\documentclass[../main.tex]{subfiles}

\begin{document}
\chapter{Basic Topology}

\section{度量空间}

\begin{definition}{}{}
    令 \(  X \neq \varnothing  \)是一个非空集合.  \(  X  \)上的一个度量, 是指一个映射 \[
    \begin{aligned}
    d: X\times X&\to \mathbb{R} \\ 
     \left( x,y\right)&\mapsto d\left( x,y \right)   
    \end{aligned}
    \]使得对于任意的 \(  x,y,z\in X  \), 有
    \begin{enumerate}
        \item \(  d\left( x,y \right)\ge 0   \), 且 \(  d\left( x,y \right)= 0   \iff  x= y\)  
        \item \(  d\left( x,y \right)= d\left( y,x \right)    \) 
        \item \(  d\left( x,z \right)\le d\left( x,y \right)+ d\left( y,z \right)     \)  
    \end{enumerate}
       
\end{definition}

\begin{example}{}{}
    \(  X\neq \varnothing  \), 定义 \(  d:X\times X\to \mathbb{R}   \) \[
    d\left( x,y \right)= \begin{cases} 0,&x= y \\ 
     1,&x\neq y\end{cases}  
    \]  是一个度量.
\end{example}
\begin{proof}
    1.2.是显然的. 3.当 \(  x= z  \)时是显然的. 若 \(  x\neq z  \), 则要么 \(  x\neq y  \), 要么 \(  z\neq y  \), 从而 \[
    d\left( x,z \right)= 1\le d\left( x,y \right)  + d\left( y,z \right) 
    \]   
\end{proof}

\section{Exercises}

\begin{problem}{Convex Bodies}{}
    A convex body is a compact, convex subset \(  K\subseteq \mathbb{R} ^{n}  \) which contains 0 as an interior point and which is symmetric around the origin .
    \begin{enumerate}
        \item Prove that the following formula defines a norm on  \(  \mathbb{R} ^{n}  \): \[
        N\left( x \right): =  \inf \left\{  \lambda > 0: \frac{x}{ \lambda }\in K \right\} 
        \] 
        \item Prove that \(  K  \) is homeomorphic to  \(  \mathbb{D}^{n}  \) and that \(   \partial K  \) is homeomorphicc to \(  \mathbb{S}^{n-1}  \)    . 
    \end{enumerate}
     
\end{problem}
\begin{note}
    \begin{enumerate}
        \item 若 \(  K  \)是一个单位球. 则 \(  N\left( x \right)= \left\| x \right\|   \).  
        \item 证明 \(  N\left( x \right)   \)是一个norm, 需要证明:1.非负正定 2.绝对一次齐次 3.三角不等式.    
        \item 由于 \(   \varphi   \)的形式足够简单, 可以通过直接构造逆映射证明双射.  
    \end{enumerate}
    
\end{note}

\begin{lemma}{}{}
    \(  X   \) is compact, \(  Y   \) is Hausdorff   . If \(  \exists f :X\to Y  \) is a continious bijection, then \(  X,Y  \) are homeomorphic.   
\end{lemma}
\begin{proof}
    \begin{enumerate}
        \item It is obvious that \(  N\left( x \right)\ge 0   \).  If \(  N\left( x \right)= 0   \) , then \[
        \inf \left\{  \lambda > 0: \frac{x }{ \lambda  }\in K  \right\}= 0
        \]   which means that for all  \(   \varepsilon > 0  \),   \(  \frac{x }{  \varepsilon   } \in K   \) .Since \(  K  \) is compact, \(  \mathrm{dist}\left( K,0 \right)< \infty   \).      If \(  x \neq 0  \) ,then \(  \left\| x \right\|> 0  \) . We take   \(   \varepsilon _0   < \frac{1 }{\mathrm{dist}\left( K,0 \right) \cdot \left\| x \right\| } \) . Then \(  \left\| \frac{x }{ \varepsilon _0  }  \right\|>   \mathrm{dist}\left( K,0 \right) \)    . But \(  \frac{x }{  \varepsilon _0  }    \in K\), which is a contradiction.For the second part, let \(  x \in \mathbb{R} ^{n}, k \in \mathbb{R} \setminus \left\{ 0 \right\}  \). For all \(   \lambda > N\left( x \right)   \) ,  \[
        \frac{x }{ \lambda  }\in K \implies  \frac{kx }{k \lambda  } \in K  
        \]  ,which shows that \(  N\left( kx \right)\le kN\left( x \right)    \)    .  Replace \(  k   \) by \(  \frac{1 }{ k}   \) and replace \(  x  \) by \(  kx  \), we have \(  N\left( kx \right)\ge kN\left( x \right)    \). Thus \(  N\left( kx \right)= kN\left( x \right)    \).       For the triangle inequality , we need to show that \[
       \inf \left\{  \lambda > 0:\frac{x_1+ x_2 }{ \lambda  }\in K  \right\} \le \inf \left\{  \lambda > 0: \frac{x_1 }{ \lambda  }\in K  \right\}+ \inf \left\{  \lambda > 0: \frac{x_2 }{ \lambda  }  \in K\right\}
        \] Let \(   \lambda _0 = N\left( x_1+ x_2 \right)   \) .  Only need to show that \[
        \frac{x_1+ x_2 }{N\left( x_1 \right)+ N\left( x_2 \right)   } \in K 
        \]For all \(   \delta > 0  \), we have  \[
        \frac{x_1 }{N\left( x_1 \right)+  \delta   }\in K,\quad \frac{x_2 }{N\left( x_2 \right)+  \delta   }  \in K
        \]. Since \(  K  \) is convex , we have the convex combination \[
        \frac{1 }{N\left( x_1 \right)+ N\left( x_2 \right)+ 2 \delta    } \left( \left( N\left( x_2 \right)+  \delta   \right)\frac{x_2 }{N\left( x_2 \right)+  \delta   }+ \left( N\left( x_1 \right)+  \delta   \right)\frac{x_1 }{N\left( x_1 \right)+  \delta   }     \right) \in K 
        \]That is \[
        \frac{x_1+ x_2 }{N\left( x_1 \right)+ N\left( x_2 \right)+ 2 \delta    }\in K 
        \] for all \(   \delta > 0  \). Then \[
        N\left( x_1+ x_2 \right)\le N\left( x_1 \right)+ N\left( x_2 \right)+ 2 \delta ,\quad \forall  \delta > 0 \implies N\left( x_1+ x_2 \right)\le N\left( x_1 \right)+ N\left( x_2 \right)      
        \] 
        \item    \(  x\in K^{\circ }\iff  N\left( x \right)< 1   \) ,  \(  x \in \mathrm{Ext}\left( K \right)\iff  N\left( x \right)> 1    \) ,  \(  x \in  \partial K  \iff N\left( x \right)= 1 \).    We define a map \(   \varphi : \mathbb{R} ^{n}\to \mathbb{R} ^{n}  \) . \[
         \varphi \left( x \right): =\begin{cases}  N\left( x \right) \frac{x }{\left\| x \right\| }  , &  x\neq 0\\ 0,& x = 0   \end{cases} 
        \]We have  \[
        \left\|  \varphi \left( x \right)  \right\|= N\left( x \right)
        \]   Then \(   \varphi   \) maps  \(  K^{\circ}  \) on to  \(  \left( \mathbb{D} ^{n}\right)^{\circ}    \), and maps \(   \partial K  \) on to \(  \mathbb{S}^{n-1}  \).  To show \(   \varphi   \) is continious , only need to show it is continious at \(  0  \) .  In fact,  \(\left\|  \lim_{x\to 0} \varphi \left( x \right)  \right\|=   \lim_{x\to 0}\left\|  \varphi \left( x \right)  \right\|=\lim_{x\to 0} N\left( x \right)= 0   \) since \(  N  \) is a norm   . We define  \(  \psi : \mathbb{R} ^{n}\to \mathbb{R} ^{n}  \) : \[
        \psi \left( x \right):= \begin{cases} \frac{\left\| x \right\| }{N\left( x \right)  }x ,& x \neq 0\\ 
         0,& x= 0  \end{cases}  
        \]   Then \(   \varphi \circ \psi = \psi \circ  \varphi = \operatorname{Id}  \) .   The same skills can be used to show \(  \psi   \) is continious at \(  0  \)  as well . 
        Finally , by the lemma above, \(  X,Y  \) are homeomorphic.  
    \end{enumerate}
    
\end{proof}

\begin{problem}{}{}
    Let \(  X,X^{\prime}   \) be topological spaces and \(  x \in X  \) , \(  x^{\prime}  \in X^{\prime}   \) be base points. Prove that the wedge sum \(  X \vee X^{\prime}   \) is homeomorphic to the subspace: \[
    \left( X\times \left\{ x^{\prime}  \right\} \right)\cup \left( \left\{ x \right\}\times X^{\prime}  \right)\subseteq  X\times X^{\prime}   
    \]    
\end{problem}

\begin{proof}
    \(  X\vee X^{\prime} = \left( X \coprod  X^{\prime}  \right)/ \left\{ x\sim x^{\prime}  \right\}   \) . Considering map \(   \varphi : X\times X^{\prime}  \to X\vee  X^{\prime}   \) \[
     \varphi \left( x_1,x_2 \right)=\begin{cases} x_1,& x_2= x^{\prime} \\ 
      x_2,&x_1= x \end{cases} 
    \] 
\end{proof}

\begin{problem}{}{}
    The torus  \(  \mathbb{T}  \)  is the quotient of \(  \left[ 0,1 \right]^{2}   \)  by the equivalence relation generated by \(  \left( 0,y \right)\sim \left( 1,y \right)    \) for all \(  x,y \in \left[ 0,1 \right]   \) . Prove that \(  \mathbb{T}  \) is homeomorphic to : 
    \begin{enumerate}
        \item The product \(  \mathbb{S}^{1}\times \mathbb{S}^{1}  \) 
        \item The quotient of \(  \mathbb{R} ^{2}  \)  under the action of the (discrete) group \(  \mathbb{Z}^{2}  \) by translation. 
    \end{enumerate}
      
\end{problem}


\begin{problem}{Projective spaces}{}
    Let \(  \mathbb{K}= \mathbb{R}  \) or \(  \mathbb{C}  \) . The \(   n  \) -dimensional projectinve space \(  \mathbb{KP}^{n}  \) is the orbit space of \(  \mathbb{k}^{n+ 1}\setminus \left\{ 0 \right\}   \) under the action of \(  \mathbb{k}^{*}  \) by rescaling. 
    \begin{enumerate}
        \item Prove that \(  \mathbb{RP}^{n}  \) is homeomorphic to the orbit space \(  \mathbb{S}^{n} \setminus \left\{ \pm 1 \right\}  \) , where \(  \left\{ \pm 1 \right\}   \) acts by multiplication on \(  \mathbb{S}^{n}\subseteq \mathbb{R} ^{n+ 1}  \) . 
        \item Prove that \(  \mathbb{CP}^{n}  \) is homeomorphic to the orbit space \(  \mathbb{S}^{2n+ 1} \setminus  \mathbb{S}^{1}  \) , where \(  \mathbb{S}^{1}= \left\{ z \in \mathbb{C}  : \left| z \right|= 1 \right\}  \) acts on \(  \mathbb{S}^{2n+ 1} \subseteq \mathbb{C}^{n+ 1}  \) by multiplication. 
        \item Prove that \(  \mathbb{RP}^{1}  \) is homeomorphic to \(  \mathbb{S}^{1}  \) and that \(  \mathbb{CP}^{1}  \) is homeomorphic to \(  \mathbb{S}^{2}  \) .             
    \end{enumerate}
           
\end{problem}
\begin{proof}
    \begin{enumerate}
        \item Define \[
         \pi: \mathbb{R} ^{n+ 1}\setminus \left\{ 0 \right\}\to \mathbb{RP}^{n} ,\quad  \pi\left( x \right)= \left[ x \right]  
        \] and define \[
        f: \mathbb{S}^{n}\to \mathbb{RP}^{n}=   \left.  \pi \right|_{\mathbb{S}^{n}}
        \] Then \(  f  \) is continuous.  Take any \(  p \in \mathbb{RP}^{n}  \), there exists \(  v \in \mathbb{R} ^{n+ 1}\setminus \left\{ 0 \right\}  \), such that \(  p= \left[ v \right]   \).  Let \(  x = \frac{v }{\left\| v \right\| }   \), then \(  x \in \mathbb{S}^{n}  \)     , \[
        f\left( x \right)= \left[ x \right]= \left[ v \right]= p   
        \] thus \(  f  \) is surjective.  If \[
        f\left( x \right)= f\left( y \right)  
        \] , then \(  \left[ x \right]= \left[ y \right]    \), there exists \(   \lambda  \in \mathbb{R}^{*}  \), such taht \(  x=  \lambda y  \). Since \[
        \left\| x \right\|= \left\|  \lambda y \right\|= \left|  \lambda  \right|\left\| y \right\|\implies \left|  \lambda  \right|= 1  
        \], \(   \lambda = \pm 1  \), \(  x= y  \) or \(  x = -y  \).Thus \(  \mathrm{Fiber}_{p}\left( f \right)= \mathrm{Orbit}_{x}\left( \left\{ \pm 1 \right\} \right)    \)  The universal property of quotient topology  gives a continuous bijection \[
        \tilde{f}: \mathbb{S}^{n}\setminus \left\{ \pm 1 \right\}\to \mathbb{RP}^{n}
        \]    Since \(  \mathbb{S}^{n}\setminus \left\{ \pm 1 \right\}  \) is compact, \(  \mathbb{RP}^{n}  \) is Hausdorff,  \(  \tilde{f}  \) is a homeomorphism.    
        \item Similarly, let \(  g:\mathbb{S}^{2n+ 1}\to \mathbb{CP}^{n}  \), \(  g\left( z \right)= \left[ z \right]    \). Take any \(  p = \left[ w \right] \in \mathbb{CP}^{n}  \) ,where \(  w \in \mathbb{C}^{n+ 1}\setminus \left\{ 0 \right\}  \). Let \(  z = \frac{w }{\left\| w \right\| }   \). Then \[
        g\left( z \right)= \left[ z \right]= \left[ w \right]= p   
        \] \(  g  \) is surjective.      
        If \[
        g\left( z_1 \right)= g\left( z_2 \right)  
        \] , then there exists \(   \lambda  \in \mathbb{C}^{*}  \), such that \(  z_1=  \lambda z_2  \), then \[
        \left\| z_1 \right\|= \left|  \lambda  \right|\left\| z_2 \right\|\implies \left|  \lambda  \right|= 1  
        \] which gives that \(   \lambda \in \mathbb{S}^{1}  \). Thus \(  \mathrm{Fiber}_{p}\left( f \right)\simeq \mathrm{Orbit}_{x}\left( \mathbb{S}^{1} \right)    \).  Thus \[
        \mathbb{CP}^{n}\simeq  \mathbb{S}^{2n+ 1}/ \mathbb{S}^{1}
        \]    
        \item \begin{enumerate}
            \item Only need to shwo that \(  \mathbb{S}^{1}/\left\{ \pm 1 \right\}\simeq \mathbb{S}^{1}  \), \[
            h\left( z \right)= z^{2} 
            \] induces a continuous bijection \(  \tilde{h}: \mathbb{S}^{1}\setminus \left\{\pm 1 \right\}\to \mathbb{S}^{1}  \). \(  \mathbb{S}^{1}\setminus \left\{ \pm 1 \right\}  \) is compact, \(  \mathbb{S}^{1}  \) is Hausdorff, \(  \tilde{h}  \) is a homemophism.
            \item Since \(  \mathbb{S}^{3}/\mathbb{S}^{1}\simeq \mathbb{S}^{2}  \).        
        \end{enumerate}
        
    \end{enumerate}
    
\end{proof}
\begin{problem}{}{}
    
The torus $\mathbb{T}$ is the quotient of $\left[ 0,1 \right] ^2$ by the equivalence relation generated by $(x,0) \sim (x,1)$ and $(0,y) \sim (1,y)$ for all $x, y \in \left[ 0,1 \right]^{2} $. Prove that $\mathbb{T}$ is homeomorphic to:
\begin{enumerate}
    \item The product $\mathbb{S}^1 \times \mathbb{S}^1$.
    \item The quotient of $\mathbb{R}^2$ under the action of the (discrete) group $\mathbb{Z}^2$ by translations.
\end{enumerate}
\end{problem}
\begin{proof}
    \begin{enumerate}
        \item Let \[
        f\left( x,y \right):=  \left( e^{2 \pi xi}, e^{2 \pi yi} \right)  
        \] Then \[
        f\left( x,0 \right)= \left( e^{2 \pi i x},1 \right)= f\left( x,1 \right)   
        \] \[
        f\left( 0,y \right)= \left( 1, e^{2 \pi yi} \right)= f\left( 1,y \right)   
        \]Thus \(  f  \) induces a map \(  \bar{f}  \) on \(  \mathbb{T}  \) such that \(  \bar{f}\left( [\left( x,y \right) ] \right)=\left( e^{2 \pi xi} , e^{2 \pi yi}\right)     \) , which is a continuous function. It is not hard to see that \(  \bar{f}: \mathbb{T}\to \mathbb{S}^{1}\times \mathbb{S}^{1}  \) is a bijection.  As a continuous image of \(  \left[ 0,1 \right]^{2}   \), \(  \mathbb{T}  \) is compact.   Take any closed subset \(  V  \) of \(  \mathbb{T}  \), \(  V  \) is compact. \(  \bar{f}\left( V \right)    \) is compact as the continuous   image of \(  V  \). Since \(  \mathbb{S}^{1}\times \mathbb{S}^{1}  \) is Hausdorff, \(  \bar{f}\left( V \right)   \) is closed.   Thus \(  \bar{f}  \) is a continuous and closed bijection, thus is a homeomorphism. 
        \item Let \[
        f\left( x,y \right)= \left( x,y \right)/ \mathbb{Z}^{2} 
        \] Then Since \(  \left( x,0 \right)-\left( x,1 \right)\in \mathbb{Z} ^{2}, \left( 0,y \right)-\left( 1,y \right)\in \mathbb{Z} ^{2}      \). \(  f  \) induces a continuous map \(  \bar{f}  \)  on from \(  \mathbb{T}  \) to \(  \mathbb{S}^{1}\times \mathbb{S}^{1}  \)   .   
        It is clear that \(  f  \) maps onto \(  \mathbb{R} ^{2}/\mathbb{Z}^{2}  \). If \(  f\left( x_1,y_1 \right)= f\left( x_2,y_2 \right)    \) , then \[
        \left( x_1-x_2,y_1-y_2 \right)\in \mathbb{Z} ^{2} \implies  x_1= x_2 \text{ or } \left\{ x_1,x_2 \right\}= \left\{ 0,1 \right\}, \text{and }y_1= y_2\text{ or } \left\{ y_1,y_2 \right\}= \left\{ 0,1 \right\}
        \]  thus \(  \left( x_1,y_1 \right)\sim \left( x_2,y_2 \right)    \) .  Again since \(  \mathbb{T}  \) is compact, \(  \mathbb{R} ^{2}/ \mathbb{Z} ^{2}  \) is Hausdorff, \(  \bar{f}  \) is a homemorphism.      
    \end{enumerate}
    
\end{proof}

\begin{problem}{}{}
    Consider the orthogonal group $O_{n-1}(\mathbb{R})$ as the subgroup of $O_n(\mathbb{R})$ of matrices of the form:
\[
\begin{pmatrix}
1 & 0 \\
0 & A
\end{pmatrix} \in O_n(\mathbb{R}), A \in O_{n-1}(\mathbb{R}).
\]
Let $O_{n-1}(\mathbb{R})$ act on $O_n(\mathbb{R})$ by left multiplication. Prove that the orbit space $O_n(\mathbb{R})/O_{n-1}(\mathbb{R})$ is homeomorphic to $\mathbb{S}^{n-1}$.
\end{problem}
\begin{proofsketch}
    将 \(  O_{n}\left( \mathbb{R}  \right)   \)刻画为 \(  \mathbb{S}^{n-1}  \)上的群作用. 此时 \(  A \sim \begin{pmatrix} 
        1&0\\ 
         0&A 
    \end{pmatrix}   \)就是北极点的稳定子群. 一但 作用是传递的, 由流形上的轨道-稳定化子定理, 就有   \[
    O_{n}\left( \mathbb{R}  \right)/ O_{n-1}\left( \mathbb{R}  \right)\simeq \mathbb{S}^{n-1}  
    \]
\end{proofsketch}
\end{document}
